%%%%%%%%%%%%%%%%%%%%%%%%%%%%%%%%%%%%%%%%%%%%%%
% Documento creado por Consuelo Quinto
% LANGEBIO - CINVESTAV
% Abril 2019

% Plantilla sencilla para la tesis de maestría y doctorado de los posgrados de Biotecnologia de Plantas y Biología Integrativa
% en CINVESTAV Unidad Irapuato y LANGEBIO
%%%%%%%%%%%%%%%%%%%%%%%%%%%%%%%%%%%%%%%%%%%%%%%
\documentclass[12pt]{article}
\usepackage[utf8]{inputenc}
\usepackage{graphicx}
\usepackage{tgtermes}
\usepackage[spanish]{babel}
\usepackage[round]{natbib}

%%%%%%%%%%%%%%%%%%%%%%%%%%%%%%%%%%%%%%%%%%%%%%%%%%%%

\begin{document}

%%% En esta sección, cambiamos el comportamiento de estos comandos:
%%% Cambiamos el nombre de la lista de tablas
\renewcommand{\listtablename}{\'Indice de tablas}
%%% Cambiamos el nombre de las tablas de "Cuadro" a "Tabla"
\renewcommand{\tablename}{Tabla} 
%%% Con este comando, quitamos el título de "Referencias" cuando insertamos la bibliografía
\renewcommand{\bibsection}{}

%%%%%%%%%%%%%%%%%%
%%% Portada
%%% Esta portada contiene el logo del CINVESTAV y los datos del tesista y su proyecto
%%% Como no queremos que haya numeracion en la primera pagina, usamos el siguiente comando para remover el estilo:
\thispagestyle{empty}
%%% Es sólo va a afectar esta pagina, la numeracion continua en la pagina 2  

%%% Usamos una tabla y una minipágina para poner el logo del CINVESTAV y el nombre de la institución 
\begin{table}[ht]
\begin{tabular}{ll}
\begin{minipage}{0.3\textwidth}
\includegraphics[scale=0.4]{Cinvestav_logo.png}
\end{minipage}
& 
\begin{minipage}{0.8\textwidth}
\begin{center}
\large{\textbf{Centro de Investigación y de Estudios
Avanzados del Instituto Politécnico Nacional \\ Unidad Irapuato} } 
\end{center}
\end{minipage}
\end{tabular}
\end{table}

%%% el comando \bigskip introduce espacios grandes entre parrafos
\bigskip
\bigskip
\bigskip
\bigskip

%%% Titulo de la tesis
\begin{center}
   \Large{\textbf{El título de tu tesis va a aquí}}
\end{center}

\bigskip
\bigskip

\begin{center}
    Tesis que presenta
\end{center}

%%% Nombre del tesista
\begin{center}
   \Large{\textbf{Aquí va tu nombre}}
\end{center}

\bigskip

%%% Grado
\begin{center}
    Para obtener el grado de \\
    \Large{Maestra en Ciencias}
\end{center}

\bigskip

%%% Nombre del posgrado
\begin{center}
    En la Especialidad de Biotecnología de Plantas
\end{center}

\bigskip
\begin{center}
% Directores de tesis
    Director de Tesis: \\
    Dr Pollito
\end{center}

\bigskip
Irapuato, Guanajuato 
\hfill
Abril 2019


%% Usamos este comando para introducir una nuev página y que todo el contenido que sigue se escriba en la siguiente página
\newpage
%%%%%%%%%%%%%%%%%%

%% Hoja que indique el lugar o lugares donde se llevó a cabo el trabajo experimental.
El presente trabajo fue realizado en ....

\newpage

%%% En esta sección escribe los agradecimientos
\textbf{Agradecimientos}\\

Gracias a todos!!

\newpage

%%% Crea la tabla de contenidos general
\tableofcontents
\newpage

%%% Crea el índice de figuras y lo agrega a la tabla de contenidos
\addcontentsline{toc}{section}{\listfigurename}
\listoffigures
\newpage

%%% Crea el índice de tablas y lo agrega a la tabla de contenidos
\addcontentsline{toc}{section}{\listtablename}
\listoftables
\newpage

%%% Usamos el asterisco para que no aparezca el numero de esta sección en la tabla de contenido
\section*{Abreviaturas} 
%%% Usamos este comando agregar esta sección a la tabla de contenidos
\addcontentsline{toc}{section}{Abreviaturas (Opcional)}


%%% Cada una de las siguientes secciones pueden ser escritas directamente en este documento
%%% O se pueden incluir como archivos separados con el comando
%%% \input{archivo.tex}

\newpage
\section*{Resumen}
\addcontentsline{toc}{section}{Resumen}

\newpage
\section*{Abstract}
\addcontentsline{toc}{section}{Abstract}

\newpage
\section{Introducción}
\input{introducion.tex}

\newpage
\section{Objetivos}

\newpage
\section{Materiales y Métodos}
\input{metodos.tex}

\newpage
\section{Resultados}


\newpage
\section{Discusión}


\newpage
\section{Perspectivas}


\newpage
\section{Referencias}
%%% Puedes modificar el estilo de bibliografía
\bibliography{referencias.bib}
\bibliographystyle{abbrvnat}
%%% Usamos este comando para que salga la lista total de las referencias
\nocite{*}

\newpage
\section*{Apéndices}
\addcontentsline{toc}{section}{Apéndices (Opcional)}


\end{document}
